% EXECUTIVE SUMMARY
\section{Executive Summary}\label{ExecutiveSummary}
\begin{spacing}{1.75}
	
	\MUFGfont{
		This validation report assesses the adequacy of \emph{\textcolor{NavyBlue}{\ModelName}} for use within the PFE framework.
		The review focuses on the FX \textbf{simulation}  and the FX \textbf{calibration}
		methodology, including conceptual soundness, data integrity, implementation considerations,
		and performance testing requirements.
	}
	
	\MUFGfont{
		\textbf{Model overview:} FX spot rates are modelled as geometric Brownian motions (GBM) under the simulation measure (Historical measure) used for exposure generation.
		USD is treated as the reporting currency: FX rates for all currencies versus USD are simulated explicitly, and cross rates are derived deterministically.
		The drift (mean function) is set by the FX forward curve implied by the domestic and foreign discount factors at the initial date.
		Volatility is calibrated from historical 5-day log returns using a three-year time series.
	}
	
	\MUFGfont{
		\textbf{Key validation themes.}
		\begin{itemize}
			\item \textbf{Conceptual soundness:} consistency of the GBM + forward-implied drift with arbitrage-free FX relations and the chosen measure;
			appropriateness of assuming lognormal spot dynamics and time-homogeneous volatility for PFE horizons.
			\item \textbf{Calibration and data:} robustness of historical vol estimation (5-day returns, 3-year window), data cleansing, missing data rules,
			and sensitivity to window length and return frequency.
			\item \textbf{Correlation and aggregation:} correctness of correlation modelling at the driver level; numerical stability and PSD corrections
			(if applied in the broader simulation engine).
			\item \textbf{Performance:} unconditional distributional tests on PIT / standardized innovations, backtesting of realized vs simulated quantiles,
			and conservativeness diagnostics relevant for PFE usage.
		\end{itemize}
	}
	
	\MUFGfont{
		\textbf{Summary assessment (high level).} The model is conceptually consistent with a standard PFE-style risk factor simulation approach for FX,
		insofar as it uses forward-implied drift and a lognormal spot process.
		However, the approach introduces limitations typical of constant-vol GBM:
		(i) inability to reproduce volatility clustering / smiles,
		(ii) potential misrepresentation under stress if calibrated vol is stale,
		(iii) potential underestimation of tail risk for long horizons or during regime shifts.
		The validation therefore emphasizes sensitivity testing, stability monitoring, and clear model use limitations.
	}
	
	\MUFGfont{
		\textbf{Provisional model risk rating:} \textbf{Amber} (subject to completion of the empirical tests described in Sections
		\ref{Performance} and \ref{DataIntegrity} and remediation / monitoring commitments).
	}
	
\end{spacing}
\clearpage
